% Task 1 — County-to-County Freight Flow Matrix
% Northeast Megaregion Freight Network Design
% All tonnage figures are in thousand short tons (ktons) unless otherwise noted.

\section{Task 1 — County-to-County Freight Flow Matrix}

The first task establishes the demand foundation for the entire network design pipeline.
A county-to-county origin--destination (O-D) freight flow matrix is constructed for the
14-state Northeast megaregion using Federal Highway Administration Freight Analysis
Framework (FAF) experimental county-level flow data, covering 2025 base-year estimates
and 2030 forecasts across all modes and commodity groups.

\subsection{Data Scope and Flow Classification}

The FAF county-county parquet file contains all-mode freight O-D flows for every pair
of US county-equivalent units. After loading and standardizing five-digit FIPS codes,
each flow record is joined against a county FIPS lookup to determine whether the origin
and destination fall within the 14-state Northeast study area (Connecticut, Delaware,
Maine, Maryland, Massachusetts, New Hampshire, New Jersey, New York, Pennsylvania,
Rhode Island, Vermont, Virginia, Washington D.C., and West Virginia — 434
county-equivalent units in total).

Every flow record is assigned to exactly one of four mutually exclusive classes:

\begin{align}
\text{flow\_type} = \begin{cases}
\textit{internal}  & \text{if both origin and destination are NE counties,}\\
\textit{inbound}   & \text{if the origin is non-NE and the destination is NE,}\\
\textit{outbound}  & \text{if the origin is NE and the destination is non-NE,}\\
\textit{outside}   & \text{otherwise (both endpoints external to the NE).}
\end{cases}
\end{align}

Only \textit{internal}, \textit{inbound}, and \textit{outbound} records are retained for
downstream analysis; \textit{outside} flows are discarded. The resulting dataset contains
33.4 million O-D-commodity records spanning 3{,}144 unique county origins and 3{,}144
unique county destinations.

\subsection{Flow-Type and Mode Breakdown}

Table~\ref{tab:task1-flow-type} summarizes total freight tonnage by flow type for
2025 and 2030. Internal flows — freight that both originates and terminates within the
Northeast — constitute the dominant share at 79.3\% of the regional freight mass.
Inbound and outbound flows each account for approximately 10\%, reflecting the
megaregion's substantial self-sufficiency combined with meaningful external trade.
All three flow types are projected to grow through 2030, with inbound freight growing
fastest (+11.6\%).

\begin{table}[H]
\centering
\caption{Freight tonnage by flow type, 2025 and 2030 (thousand short tons).}
\begin{tabular}{lrrr}
\toprule
Flow Type & 2025 (ktons) & 2030 (ktons) & Growth (\%) \\
\midrule
Internal  & 2{,}624{,}252 & 2{,}879{,}815 & +9.7 \\
Inbound   &   344{,}774 &   384{,}860 & +11.6 \\
Outbound  &   338{,}619 &   360{,}668 & +6.5 \\
\midrule
\textbf{Total} & \textbf{3{,}307{,}644} & \textbf{3{,}625{,}342} & +9.6 \\
\bottomrule
\end{tabular}
\label{tab:task1-flow-type}
\end{table}

Table~\ref{tab:task1-mode} presents the modal breakdown. Truck-and-air (mode 11)
is overwhelmingly dominant at 72.1\% of total NE freight tonnage, confirming that
road-based freight logistics is the primary planning focus. Pipeline freight (mode 6)
accounts for 15.3\% — largely petroleum and natural gas — and is expected to grow only
marginally (+0.6\%) through 2030. Rail contributes 6.1\%, while water freight is
projected to decline slightly ($-$4.6\%), reflecting modal shifts on the Eastern Seaboard.

\begin{table}[H]
\centering
\caption{Freight tonnage by transport mode, 2025 and 2030 (thousand short tons).}
\begin{tabular}{llrrr}
\toprule
Mode & Label & 2025 (ktons) & 2030 (ktons) & Growth (\%) \\
\midrule
11 & Truck + Air        & 2{,}384{,}432 & 2{,}680{,}331 & +12.4 \\
6  & Pipeline           &   505{,}769 &   508{,}568 & +0.6  \\
2  & Rail               &   201{,}009 &   209{,}873 & +4.4  \\
5  & Multiple / Mail    &   125{,}227 &   139{,}551 & +11.4 \\
3  & Water              &    91{,}207 &    87{,}019 & $-$4.6 \\
\bottomrule
\end{tabular}
\label{tab:task1-mode}
\end{table}

The trade-type distribution (Table~\ref{tab:task1-trade}) shows that 87.5\% of
NE freight is purely domestic in character (trade type 1 — internal domestic).
International import flows (6.6\%) and export flows (5.9\%) are growing faster
than domestic freight through 2030, indicating increasing integration with global
supply chains.

\begin{table}[H]
\centering
\caption{Freight tonnage by trade type, 2025 and 2030 (thousand short tons).}
\begin{tabular}{llrrr}
\toprule
Trade Type & Label & 2025 (ktons) & 2030 (ktons) & Growth (\%) \\
\midrule
1 & Internal (domestic) & 2{,}894{,}433 & 3{,}157{,}031 & +9.1  \\
2 & Import              &   217{,}705 &   249{,}241 & +14.5 \\
3 & Export              &   195{,}506 &   219{,}069 & +12.1 \\
\bottomrule
\end{tabular}
\label{tab:task1-trade}
\end{table}

\subsection{State-Level Summaries and Pareto Analysis}

Table~\ref{tab:task1-state} reports per-state inbound and outbound tonnage for all 14
NE states in 2025. Pennsylvania leads in outbound freight (694{,}901 ktons), driven by
its role as a regional distribution hub and gateway between the Midwest and the
Northeast coast. New York exhibits the largest inbound volume (612{,}371 ktons),
reflecting the gravitational pull of the New York metropolitan area. Virginia and
New Jersey together account for over 725{,}000 ktons of outbound freight, underpinned
by proximity to major port and intermodal complexes. The smallest states — Rhode Island,
Vermont, and Washington D.C. — generate less than 50{,}000 ktons each in outbound flows.

\begin{table}[H]
\centering
\caption{Per-state freight summary, 2025 (thousand short tons).}
\begin{tabular}{lrr}
\toprule
State & Outbound 2025 (ktons) & Inbound 2025 (ktons) \\
\midrule
Pennsylvania      & 694{,}901 & 597{,}803 \\
New York          & 534{,}471 & 612{,}371 \\
Virginia          & 371{,}909 & 420{,}584 \\
New Jersey        & 353{,}849 & 382{,}218 \\
West Virginia     & 230{,}859 & 139{,}038 \\
Maryland          & 214{,}835 & 233{,}391 \\
Massachusetts     & 180{,}897 & 195{,}114 \\
Connecticut       & 125{,}714 & 132{,}590 \\
Maine             &  78{,}195 &  73{,}460 \\
New Hampshire     &  57{,}150 &  65{,}844 \\
Delaware          &  45{,}071 &  47{,}411 \\
Rhode Island      &  36{,}589 &  35{,}068 \\
Vermont           &  32{,}092 &  24{,}332 \\
Washington D.C.   &   6{,}338 &   9{,}801 \\
\bottomrule
\end{tabular}
\label{tab:task1-state}
\end{table}

A Pareto analysis of county-level outbound truck tonnage identifies the top-50 origin
counties responsible for the largest outbound freight volumes. Table~\ref{tab:task1-top10}
lists the top ten. Allegheny County (Pittsburgh, PA) leads all NE origins at
65{,}815 ktons in 2025, followed by Erie County (Buffalo, NY) at 54{,}844 ktons and
New York County (Manhattan, NY) at 52{,}158 ktons. The top-50 origin counties are
concentrated in the major metropolitan corridors: Greater Philadelphia--Pittsburgh
(PA), Greater New York (NY/NJ), and Greater Boston (MA), consistent with the I-95
and I-76 freight corridors. On the destination side, Erie County, NY ranks first
(61{,}472 ktons) followed closely by Allegheny County, PA (60{,}339 ktons), reinforcing
the bidirectional importance of the PA--NY freight corridor.

\begin{table}[H]
\centering
\caption{Top 10 origin counties by outbound freight tonnage, 2025 (thousand short tons).}
\begin{tabular}{rlllr}
\toprule
Rank & FIPS & County & State & Outbound 2025 (ktons) \\
\midrule
1  & 42003 & Allegheny & Pennsylvania   & 65{,}815 \\
2  & 36029 & Erie      & New York       & 54{,}844 \\
3  & 36061 & New York  & New York       & 52{,}158 \\
4  & 42101 & Philadelphia & Pennsylvania & 45{,}514 \\
5  & 34023 & Middlesex & New Jersey     & 44{,}082 \\
6  & 25017 & Middlesex & Massachusetts  & 42{,}700 \\
7  & 24031 & Montgomery & Maryland      & 41{,}319 \\
8  & 34039 & Union     & New Jersey     & 39{,}042 \\
9  & 42123 & Warren    & Pennsylvania   & 36{,}845 \\
10 & 36067 & Onondaga  & New York       & 36{,}439 \\
\bottomrule
\end{tabular}
\label{tab:task1-top10}
\end{table}

The constructed O-D matrix provides the empirical demand foundation for all downstream
tasks. County-level throughput values — derived from this matrix by summing inbound and
outbound flows at each county endpoint — serve as the demand weights for region
clustering (Task 3) and the capacity sizing inputs for hub location (Task 5). The
dominance of truck-compatible freight confirms that the regional hub network will be
designed primarily around road access, with rail and intermodal connectivity serving
as secondary criteria.
